\section{Conclusion}

This thesis investigated perturbation-based learning as a biologically motivated alternative to backpropagation. The work was organized around three research questions, which are answered directly below.

\textbf{What has been proposed and theoretically established in the literature on\\ perturbation-based learning?}  
The thesis concludes that perturbation-based learning is an established but still incomplete research direction. In Chapter~2, we identified 21 different update rules across 9 articles and organized them in a common notation. This showed that there exist many different versions of perturbation-based learning, with variants that differ in where noise is injected, how the learning signal is formed, and how the update is scaled. The two central families of perturbation methods are weight perturbation, where noise is applied directly to the parameters, and node perturbation, where noise is applied to neural activity or pre-activations. However, despite the existence of many relevant articles and update rules, the relationship between WP and NP remains insufficiently resolved in the current literature. While previous work provides an analytical comparison showing why NP can be expected to outperform WP in a simplified one-layer linear setting, no corresponding theoretical derivation appears to exist for how the two families compare in general nonlinear multilayer networks. The conclusion is therefore that perturbation-based learning is well established as a family of learning rules, but that its theoretical foundation remains incomplete because the relationship between its two central families has not been fully clarified.

\textbf{How can perturbation-based learning be further developed?}  
This is the most important question of the thesis. The answer is that perturbation-based learning can be further developed by clarifying the relationship between WP and NP, which was identified above as an incomplete part of the existing literature. Inspired by the local reparameterization trick of \textcite{kingma_variational_2015}, Chapter~3 showed that the weight-level noise in WP can be translated into an equivalent pre-activation noise distribution. This led to the derivation of input-scaled node perturbation (IS-NP), a node-level update rule with the same expected update as WP, but with a variance that is theoretically guaranteed not to be larger and is expected to be lower in practice. The empirical results support this prediction: WP has much higher estimator variance than IS-NP, and performs clearly worse than the NP methods on the more difficult tasks. This suggests that the variance reduction is not only a theoretical result, but also matters for practical learning, and supports NP as the more promising perturbation-learning direction. IS-NP also performs best among the NP methods included in the evaluation, outperforming both vanilla NP and fan-in-scaled NP in the most challenging experiments. The conclusion is therefore that the thesis further develops perturbation-based learning in two connected ways: by strengthening its theoretical foundation through the WP--NP clarification, and by deriving a new perturbation method that appears practically promising in the settings studied here.

\textbf{What is the practical capability of perturbation-based learning?}  
The empirical results show that perturbation-based learning is capable of training nontrivial multilayer networks, including meaningful performance on MNIST classification using an MLP. However, the methods remain clearly limited compared with backpropagation. The gap is small on the simplest task, becomes more visible on California housing and MNIST, and is large on CIFAR-10, where the perturbation methods do not approach the speed or final performance of backpropagation. This suggests that current perturbation-based methods are not sufficient to train state-of-the-art models or solve highly complex supervised learning problems competitively on their own. Their significance is instead that they demonstrate that biologically motivated learning rules based on stochastic perturbations and scalar feedback can support genuine learning. This should not be interpreted as evidence that perturbation-based learning will necessarily become a central component of future AI systems. The experiments were carried out in standard artificial neural networks, not in biologically realistic architectures, and the broader relevance of these methods remains speculative. A more cautious conclusion is that perturbation-based learning may be one useful building block among several in future biologically inspired learning systems, much like three-factor learning is one learning mechanism among several in the brain.

Whether perturbation-based learning will ever become important for future AI systems is impossible to know. It may remain a biologically motivated curiosity, useful mainly for understanding a narrow class of learning rules. But it may also turn out that some of the next major steps in AI require ideas that look quite different from today’s dominant recipe of larger models, larger datasets, and more backpropagation. If that future includes architectures inspired more deeply by biology, then perturbation-based learning could be one of the mechanisms that makes learning in such systems possible. The contributions of this thesis make that possibility a little more plausible.